\sheet{2}{OpenMP basics, cores, and placement}{1h30}
\section{Preparation}
Define \textbf{physical core}, \textbf{logical processor}, \textbf{thread},
\textbf{parallel region}, \textbf{worksharing}, \textbf{reduction} and
\textbf{affinity}. Draw one process with several threads and indicate their
shared address space. Distinguish shared reading from conflicting writes.

Sources are in \code{02-openmp-basic}. Each file has its own
\code{main}; the implementations of other exercises are not required. B3 is
supplied in full for observation. Commands run from this \code{02-openmp-basic} directory:
\begin{lstlisting}[style=shell]
cmake -S . -B build -DCMAKE_BUILD_TYPE=Release
\end{lstlisting}
Use the allocated CPUs only. \code{OMP_NUM_THREADS} requests a team size;
\code{OMP_DYNAMIC=false} disables dynamic team-size adjustment, but other limits
may still apply. Environment assignments before a command apply to that launch.
Rebuild after edits. Each numerical program checks the requested input before reporting timings.

This folder contains its own CMake files and utility headers; it can be copied
and compiled independently. Numerical sizes are freely chosen integers within
the ranges listed in \code{README.md}; the examples are not a predefined menu.
For programs whose usage lists \code{--self-test}, add that flag to run extra
small and edge cases before the requested input. Without it, only the requested
size is checked and measured. Rebuild after changing the implementation.

\section{Exercises}
\exercise{B1}{Distributing an independent loop}
Edit \code{B1-affine.cpp}. The required operation is $y_i=2x_i+1$ on separate
input and output arrays. An ordinary loop inside a parallel region is executed
in full by every thread. Writing the same value concurrently does not remove a race.
\begin{enumerate}
\item Analyse this incorrect fragment without executing it:
\begin{lstlisting}
#pragma omp parallel
{
    for (std::size_t i = 0; i < n; ++i)
        y[i] = 2*x[i] + 1;
}
\end{lstlisting}
With four threads and eight elements, how many times is each iteration attempted?
Identify the missing construct.
\item Implement \code{affine} using a worksharing loop and
\code{default(none)}. Specify the data-sharing attributes that need clauses.
Explain why the associated loop index is private and why each output has one writer.
\item Build and run:
\begin{lstlisting}[style=shell]
cmake --build build --target b1_affine -j
OMP_NUM_THREADS=4 OMP_DYNAMIC=false ./build/bin/b1_affine 1048576 7
\end{lstlisting}
The arguments are length and trials. The supplied opening parallel region prints
one complete line per thread inside a critical region. Explain why the order of
these lines can vary. Keep printing outside arithmetic timings.
\item Repeat with 1, 2 and 4 requested threads. Add \code{--self-test} to check empty input and
uneven small lengths. For seven elements and three threads, justify correctness
without assuming equal chunks or a particular assignment.
\end{enumerate}

\exercise{B2}{Reduction, atomic updates and critical regions}
Edit \code{B2-dot.cpp}. The dot product is $d=\sum_i x_i y_i$.
The products are independent, but every contribution must enter the final sum.
\begin{enumerate}
\item Starting with zero and contributions 2 and 5, draw an interleaving of
ordinary reads and writes that loses a contribution. Explain why a barrier
after such a loop cannot repair the race.
\item Implement \code{dot_reduction} with an addition reduction. Initialise
the original result to zero. OpenMP creates private contributions and combines
them before the function returns.
\item Implement \code{dot_atomic} and \code{dot_critical} with one protected
shared update per element. Use an atomic update and a named critical region,
respectively. These versions intentionally coordinate every addition.
\item Build and run:
\begin{lstlisting}[style=shell]
cmake --build build --target b2_dot -j
OMP_NUM_THREADS=4 OMP_DYNAMIC=false ./build/bin/b2_dot 1048576 7
\end{lstlisting}
The arguments are length and trials. All three functions in this file must be
implemented. Checks compare each result with a reference; floating-point tolerance
allows rounding differences from regrouping.
\item Repeat with 1, 2 and 4 threads. Record the actual team size, median and
range for each variant. Calculate $S_p=T_1/T_p$ and $E_p=S_p/p$ using each
variant's own one-thread time. Explain why more threads can make a correct
implementation slower through coordination overhead or limited memory bandwidth.
\end{enumerate}

\exercise{B3}{Observing thread placement}
Run the supplied \code{B3-affinity.cpp}. On Linux, identify the allowed logical
CPUs and their core/socket membership:
\begin{lstlisting}[style=shell]
lscpu -e
taskset -pc $$
cmake --build build --target b3_affinity -j
\end{lstlisting}
Two logical processors may share a physical core. Use both socket and core IDs
when identifying distinct cores.
\begin{enumerate}
\item Compare placement policies:
\begin{lstlisting}[style=shell]
OMP_NUM_THREADS=4 OMP_DYNAMIC=false OMP_PLACES=cores \
  OMP_PROC_BIND=false ./build/bin/b3_affinity 100000
OMP_NUM_THREADS=4 OMP_DYNAMIC=false OMP_PLACES=cores \
  OMP_PROC_BIND=close ./build/bin/b3_affinity 100000
OMP_NUM_THREADS=4 OMP_DYNAMIC=false OMP_PLACES=cores \
  OMP_PROC_BIND=spread ./build/bin/b3_affinity 100000
\end{lstlisting}
The argument is samples per thread. Distinguish the allowed CPU set from the CPU
observed at one instant. \code{close} and \code{spread} operate on the runtime's
available places, so their mappings need not differ on every machine.
\item Choose an allowed logical CPU ID and replace \code{ID} below:
\begin{lstlisting}[style=shell]
OMP_NUM_THREADS=4 OMP_DYNAMIC=false OMP_PROC_BIND=false \
  taskset -c ID ./build/bin/b3_affinity 100000
\end{lstlisting}
Explain how four threads progress when restricted to that one logical CPU.
Repeat with two allowed IDs separated by a comma.
\item Rerun B2 using \code{OMP_PLACES=cores} and an explicit binding policy,
without the one-CPU restriction. Record placement with the timings. Why does a
correct result not prove that four distinct physical cores were used?
\item If no CPU changes were sampled, what can and cannot be concluded?
This diagnostic's own runtime includes sampling and printing; it is not a
benchmark of the dot product. On non-Linux systems, use the OpenMP runtime's
affinity display for placement observations.
\end{enumerate}

\exercise{B4}{Alignment of an address}
Edit \code{B4-alignment.cpp}. An address is aligned to $A$ bytes when it is
divisible by $A$. Being aligned to 8 bytes does not imply alignment to 64 bytes.
\begin{enumerate}
\item Implement \code{power_two_alignment} on an integer address. Return the
largest power of two dividing a nonzero address; define the result for zero as
zero. Do not dereference invented addresses.
\item Predict results for 1, 6, 8, 24 and 64. Build and run:
\begin{lstlisting}[style=shell]
cmake --build build --target b4_alignment -j
./build/bin/b4_alignment
\end{lstlisting}
The program also tests zero and a high-bit address, then prints addresses for a
64-byte-aligned array and its second element. It takes no numerical arguments.
\item Assuming 8-byte doubles, derive the second element's guaranteed alignment.
Distinguish the declaration's guarantee from a larger alignment observed by chance.
Explain why ordinary allocation respects the type's alignment but need not put
every object on a cache-line boundary.
\end{enumerate}
